\clearpage
\begin{trapbox}
\textbf{易错、边界与最小反例（C091 专项核查）。}
\begin{itemize}
\item \textbf{把面积错写成 $ab$ 或 $2ab$。}
必有系数 $\boxed{1/2}$。\textbf{最小核对：}
$a=3,b=2,P=(3,0)$，由 $Q=(3/2,1),R=(3/2,-1)$ 算出面积 $3$，不是 $6$。

\item \textbf{以为 $P$ 是某个边上的点。}
$P$ 与两渐近线交点 $O$ 是\textbf{对顶点}；
另两顶点 $Q,R$ 各自\textbf{落在不同渐近线}上。

\item \textbf{错误地把两组平行线间距的乘积当作面积。}
一般应是 $S=d_+d_-/\sin\theta$，而非 $d_+d_-$；
只有两渐近线\textbf{垂直}（$a=b$）时，$\sin\theta=1$。

\item \textbf{把一般平行四边形默认为矩形或菱形。}
仅 $a=b$ 时两渐近线垂直，得到矩形；
仅在 $P=(\pm a,0)$ 时相邻两边等长，得到菱形。
两者兼有时才是正方形。

\item \textbf{菱形、矩形条件只背结论，忽略由来。}
记 $Q=(u,ku),R=(v,-kv)$，$k=b/a$。
因 $4uv=a^2>0$，$u,v$ 同号；
$|OQ|=|u|\sqrt{1+k^2}$，$|OR|=|v|\sqrt{1+k^2}$。
故\textbf{菱形} $\Leftrightarrow|u|=|v|\Leftrightarrow u=v
\Leftrightarrow y=k(u-v)=0\Leftrightarrow P=(\pm a,0)$。
\textbf{矩形} $\Leftrightarrow$ 渐近线垂直 $\Leftrightarrow
(a,b)\cdot(a,-b)=a^2-b^2=0\Leftrightarrow a=b$。
\textbf{正方形}需两者同时满足。不要把某个$P$的菱形条件错归因于$a=b$。

\item \textbf{把渐近线写成 $y=\pm a x/b$。}
对 $x^2/a^2-y^2/b^2=1$，应写
$\boxed{y=\pm (b/a)x}$；先在方程中取最高次项 $x^2/a^2-y^2/b^2=0$ 即可恢复。

\item \textbf{认为只能在右支或 $y\ge0$ 才成立。}
左支同样成立。平方和面积的绝对值使 $x,y$ 的符号不影响结果。

\item \textbf{混淆 $a,b$ 与焦距 $c$。}
这个面积定值只涉及\textbf{两条半轴参数的乘积 $ab$}，
与 $c=\sqrt{a^2+b^2}$ 不需直接联立。

\item \textbf{误用椭圆公式。}
同取两条直线 $y=\pm(b/a)x$，若 $P$ 在
椭圆 $x^2/a^2+y^2/b^2=1$ 上，面积并不恒定：
$P=(a,0)$ 时为 $ab/2$；$P=(a/\sqrt2,b/\sqrt2)$ 时四边形退化、面积为 $0$。

\item \textbf{反向推断失败。}若仅知道 $S=ab/2$，
则 $\bigl|x^2/a^2-y^2/b^2\bigr|=1$，
有 $+1$、$-1$ \textbf{两种情形}，不能唯一推出题设双曲线方程。

\item \textbf{忘记渐近线上是退化情形。}
当 $P$ 落在 $\ell_+$ 或 $\ell_-$ 上时，
$u=0$ 或 $v=0$，所谓“四边形”有相邻点重合（或共线），
只可把它\textbf{看成退化平行四边形}，面积按 $0$ 计。
不能说它仍是通常意义的非退化平行四边形。

\item \textbf{将固定 $c$ 与固定 $a,b$ 混为一谈。}
若 $c^2=a^2+b^2$ 固定而 $a,b$ 可以变化，则
$0<S=ab/2\le c^2/4$，只有上界，面积\textbf{不固定}；
且 $S=0$ 不可取得。若 $a,b$ 固定，面积才与点 $P$ 的运动位置无关。

\item \textbf{套用一次式乘积推广却忽略独立性。}
若 $L_1,L_2$ 方向相同（系数行列式 $D=0$），
两组直线不能围成所需平行四边形；
公式 $|L_1L_2|/|D|$ 不可使用。

\item \textbf{忘记 $a,b>0$ 或面积非负。}
方程使用 $a^2,b^2$，约定 $a,b>0$；
推广 $x^2/a^2-y^2/b^2=t$ 时必须写
$\boxed{S=ab|t|/2}$ 而不是 $abt/2$。

\item \textbf{拿错两条平行线。}必须\textbf{一条平行 $\ell_+$，另一条平行 $\ell_-$}；
若两条都平行一条渐近线，无法按题意围成该平行四边形。
\end{itemize}

\textbf{考场自检：}
\[
\boxed{\ell_\pm:y=\pm\frac ba x}\quad
\boxed{P=Q+R}\quad
\boxed{S=\frac{ab}{2}}\quad
\boxed{\text{两支均可}}.
\]
\end{trapbox}

\clearpage
\begin{trapbox}
\textbf{模型对照：换曲线后，面积还恒定吗？}
\begin{center}
% Equal scale a=3,b=2. Ellipse points show area 3 vs 0;
% the third point lies on the conjugate hyperbola and also gives area 3.
\begin{tikzpicture}[x=.62cm,y=.62cm,>=stealth,font=\small]
\foreach \s/\title in {0/{椭圆：$P=(3,0)$},7.6/{椭圆：$P=(3/\sqrt2,\sqrt2)$},15.2/{共轭双曲线：$P=(0,2)$}}{
\begin{scope}[shift={(\s,0)}]
\node at (0,3.4) {\title};
\draw[gray,dashed] (-3.3,-2.2)--(3.3,2.2) (-3.3,2.2)--(3.3,-2.2);
\node[below left] at (0,0) {$O$};
\end{scope}}
\begin{scope}
\draw[teal!70!black,thick] (0,0) ellipse[x radius=3,y radius=2];
\fill[orange!17] (0,0)--(1.5,1)--(3,0)--(1.5,-1)--cycle;
\draw[orange!85!black,thick] (0,0)--(1.5,1)--(3,0)--(1.5,-1)--cycle;
\fill[red!80!black] (3,0) circle(2pt) node[right] {$P$};
\node at (0,-3) {$S=3$};
\end{scope}
\begin{scope}[shift={(7.6,0)}]
\draw[teal!70!black,thick] (0,0) ellipse[x radius=3,y radius=2];
\draw[orange!85!black,very thick] (0,0)--(2.12132,1.41421);
\fill[red!80!black] (2.12132,1.41421) circle(2pt) node[above left=7pt] {$P=Q$};
\node at (0,-3) {$R=O,\ S=0$（退化）};
\end{scope}
\begin{scope}[shift={(15.2,0)}]
\draw[teal!70!black,thick,domain=-3:3,samples=60] plot (\x,{2*sqrt(1+\x*\x/9)});
\draw[teal!70!black,thick,domain=-3:3,samples=60] plot (\x,{-2*sqrt(1+\x*\x/9)});
\fill[orange!17] (0,0)--(1.5,1)--(0,2)--(-1.5,1)--cycle;
\draw[orange!85!black,thick] (0,0)--(1.5,1)--(0,2)--(-1.5,1)--cycle;
\fill[red!80!black] (0,2) circle(2pt) node[above] {$P$};
\node at (0,-3.7) {$S=3,\quad P\notin H$};
\end{scope}
\end{tikzpicture}

\end{center}
三图统一取 $a=3,b=2$、同一比例。左、中两图在同一椭圆上，面积却从 $3$ 变为 $0$；
右图满足 $x^2/9-y^2/4=-1$，面积也为 $3$。
所以既不能把双曲线定值搬到椭圆，也不能由面积 $ab/2$ 反推原双曲线。

\medskip
\textbf{形状对照：矩形与菱形的条件各不相同。}
\begin{center}
% All panels use the same physical scale. Points are vertices of H.
\begin{tikzpicture}[x=.95cm,y=.95cm,font=\small]
\foreach \s/\aa/\bb/\tt in {0/3/2/{菱形},4.7/2/2/{正方形},9.4/2/2/{矩形}}{
\begin{scope}[shift={(\s,0)}]
\node at (1.5,2.8) {\tt};
\ifdim\s pt=9.4pt
  % a=b=2, lambda=1,mu=1/4; P=(5/2,3/2).
  \fill[orange!17] (0,0)--(2,2)--(2.5,1.5)--(.5,-.5)--cycle;
  \draw[orange!85!black,thick] (0,0)--(2,2)--(2.5,1.5)--(.5,-.5)--cycle;
  \node[right] at (2.5,1.5) {$P$};
\else
  \fill[orange!17] (0,0)--(\aa/2,\bb/2)--(\aa,0)--(\aa/2,-\bb/2)--cycle;
  \draw[orange!85!black,thick] (0,0)--(\aa/2,\bb/2)--(\aa,0)--(\aa/2,-\bb/2)--cycle;
  \node[right] at (\aa,0) {$P$};
\fi
\node[left] at (0,0) {$O$};
\node at (1.5,-1.65) {$a=\aa,\ b=\bb$};
\end{scope}}
\end{tikzpicture}

\end{center}
菱形要求 $|OQ|=|OR|$，本题等价于 $P=(\pm a,0)$；
矩形要求两渐近线垂直，即 $a=b$。两条件同时满足才是正方形。
\textbf{判定两步：} $|OQ|=|OR|\Leftrightarrow|u|=|v|$，
而 $uv>0$ 强制 $u=v$，因此 $y=0$；
另一方面 $(a,b)\cdot(a,-b)=a^2-b^2$，
正交当且仅当 $a=b$。
\end{trapbox}
